\documentclass[12pt]{article}
\usepackage{amsfonts}
\usepackage{amsmath}
\usepackage{amssymb}
\usepackage{amsthm}
\usepackage{amssymb}
\usepackage{graphicx}
\usepackage{enumerate}
% \usepackage{enumitem}
\usepackage{xfrac}
\usepackage{xcolor}


\usepackage[margin=1in]{geometry}

% \usepackage{ulem} % for the strikethrough \sout{}

\usepackage{hyperref}

% this is apparently really good?
\usepackage{mathtools}

% AI generated box for the problem statement
\usepackage[most]{tcolorbox}
\newtcolorbox{problem}[1][]{
    colback=gray!5!white,    % Background color: 5% gray, 95% white
    colframe=gray!80!black,  % Border color: 80% gray, 20% black
    fonttitle=\bfseries,     % Make the title text bold
    title={#1},              % Use the optional argument as the title
    arc=3pt,                 % Rounded corners (set to 0pt for sharp corners)
    boxrule=1pt,             % Thickness of the border
    breakable                % Allows the box to split across page boundaries
}

% this changes the footnotes to not use numbers (this is better in equations IMO,
% doesn't look like an exponent)
\usepackage[symbol*]{footmisc}
\setfnsymbol{wiley}

%bibliography stuff
\usepackage[
backend=biber,
style=numeric,
sorting=ynt
]{biblatex}
\addbibresource{bib.bib}

% I use this to make the page with the table landscape mode
\usepackage{pdflscape}

% use this to wrap text nicely in columns, with the X column type
\usepackage{tabularx}

% I'm copying the textbook's notation for these things,
% I defined them up here so I can change them as easily as possible
% \newcommand{\R}{{\mathbf{R}}}
% \newcommand{\Rn}[1][n]{{\mathbf{R^{#1}}}} % for some reason they bold the exponent, weird
\newcommand{\R}{{\mathbb{R}}}
\newcommand{\C}{{\mathbb{C}}}
\newcommand{\D}{{\mathbb{D}}}
\newcommand{\Hplane}{{\mathbb{H}}}
% \newcommand{\Rn}[1][n]{{\mathbb{R}^{#1}}}
\newcommand{\Q}{{\mathbb{Q}}}
\newcommand{\N}{{\mathbb{N}}}
\newcommand{\Z}{{\mathbb{Z}}}
\newcommand{\comp}[1]{#1^C}

% extended reals
\newcommand{\ER}{[-\infty, +\infty]}

% convergence in measure
\newcommand{\tomeasure}{\overset{m}{\to}}

% support
\newcommand{\supp}{\operatorname{supp}}

% indicator function
\newcommand{\ind}[1]{\chi_{#1}}

% you can change this from : to \middle\vert depending on style choices I guess
\newcommand{\setbuild}[2]{\left\{ {#1} \; \middle\vert \; {#2} \right\}}

\newcommand{\abs}[1]{\left\vert #1 \right\vert}
\newcommand{\norm}[1]{\left\| #1 \right\|}
\newcommand{\vect}[1]{\mathbf{{#1}}}
\newcommand{\diag}{\operatorname{diag}}

\newcommand{\measure}[1]{\left\vert #1 \right\vert}

\newcommand{\omeasure}[1]{\abs{#1}_e}

\newcommand{\closure}[1]{\overline{#1}}
\newcommand{\interior}[1]{{#1}^\circ}
\newcommand{\bdy}{\partial} % maybe I'll want to change this notation later

\newcommand{\BV}{\operatorname{BV}}
\newcommand{\AC}{\operatorname{AC}}

\newcommand{\WeakLp}[1][p]{\operatorname{Weak} L^{#1}}
\newcommand{\pnorm}[2][p]{\norm{#2}_{L^{#1}}}
\newcommand{\lpnorm}[2][p]{\norm{#2}_{l^{#1}}}

% we use these
\newcommand{\upto}{\mathrel{\nearrow}}
\newcommand{\downto}{\mathrel{\searrow}}

\newcommand{\ft}[1]{\widehat{#1}}
\newcommand{\ip}[1]{\left\langle{#1}\right\rangle}

% make \iff look better
\renewcommand{\iff}{\Leftrightarrow}
\renewcommand{\implies}{\Rightarrow}
\renewcommand{\impliedby}{\Leftarrow}

% change the default phi and epsilon
\renewcommand{\epsilon}{\varepsilon}
% \renewcommand{\phi}{\varphi}

% From 660
\renewcommand{\Re}{\operatorname{Re}}
\renewcommand{\Im}{\operatorname{Im}}
\newcommand{\cc}[1]{\overline{#1}}
\newcommand{\Res}{\operatorname{Res}}
\newcommand{\coker}{\operatorname{coker}} %why was this in there lol?


% cute
\newcommand{\RHS}{\textrm{RHS}}
\newcommand{\LHS}{\textrm{LHS}}


%%% Formatting for theorems and so on

% AI code
% Define the custom theorem style
\newtheoremstyle{ColorTheorem}% Name of the style
  {\topsep}% Space above
  {\topsep}% Space below
  {\upshape}% Body font (usually italics for theorems)
  {}% Indent amount
  {\bfseries\color{blue}}% Theorem head font (makes "Theorem 1" blue and bold)
  {.}% Punctuation after theorem head
  {5pt plus 1pt minus 1pt}% Space after theorem head
  {\thmname{#1}\thmnumber{ #2}\thmnote{ \textcolor{red}{(#3)}}}% Theorem head spec
\newtheoremstyle{ColorDefinition}% Name of the style
  {\topsep}% Space above
  {\topsep}% Space below
  {\upshape}% Body font (usually italics for theorems)
  {}% Indent amount
  {\bfseries\color{blue}}% Theorem head font (makes "Theorem 1" blue and bold)
  {.}% Punctuation after theorem head
  {5pt plus 1pt minus 1pt}% Space after theorem head
  {\thmname{#1}\thmnumber{ #2}\thmnote{ \textcolor{red}{(#3)}}}% Theorem head spec

\newtheoremstyle{ColorRemark}% Name of the style
  {\topsep}% Space above
  {\topsep}% Space below
  {\upshape}% Body font (usually italics for theorems)
  {}% Indent amount
  {\bfseries\color{black}}% Theorem head font (makes "Theorem 1" blue and bold)
  {.}% Punctuation after theorem head
  {5pt plus 1pt minus 1pt}% Space after theorem head
  {\thmname{#1}\thmnumber{ #2}\thmnote{ \textcolor{red}{(#3)}}}% Theorem head spec

\theoremstyle{ColorTheorem}
\newtheorem{theorem}{Theorem}
\providecommand*\lemmaautorefname{Theorem}
\numberwithin{theorem}{section}
\newtheorem{corollary}{Corollary}
\providecommand*\lemmaautorefname{Corollary}
\numberwithin{corollary}{section}
\newtheorem{proposition}[theorem]{Proposition}
\providecommand*\lemmaautorefname{Proposition}
\newtheorem{lemma}[theorem]{Lemma}
\providecommand*\lemmaautorefname{Lemma}

\theoremstyle{ColorDefinition}
\newtheorem{definition}{Definition}[section]
\theoremstyle{definition}

\theoremstyle{ColorRemark}
\newtheorem*{remark}{Remark}

% \renewcommand{\thesection}{(\arabic{section})}
% reset the section numbering in the second part
\makeatletter
\@addtoreset{section}{part}
\makeatother

% \author{Joshua Rippeon}
% \title{Key Information}
% \date{May 10, 2026}

% % Wheeden and Zygmund use this convention
% \renewcommand{\setminus}{-}

% use section numbers for equations
\numberwithin{equation}{section}

% use this to number a specific line in an align
\newcommand\numberthis{\addtocounter{equation}{1}\tag{\theequation}}

% allows align* and similar to cross page boundaries
\allowdisplaybreaks

% \includeonly{
%   wzch2.tex
% }

% % emulate the typesetting from the 1979 printing of Ahlfors (good vibes)
% \usepackage{mlmodern}

\title{AMSC665 Project 2}
\author{Joshua Rippeon}

\begin{document}

\maketitle


\tableofcontents
%%%%%%%%%%%%%%%%%%%%%%%%%%%%%%%%%%%%%%%%%%%%%%%%%%%%%%%%%%%%%%%%%%%%%%%%%%%%%%%%
\section{Variational neural network for the committor problem.}
%%%%%%%%%%%%%%%%%%%%%%%%%%%%%%%%%%%%%%%%%%%%%%%%%%%%%%%%%%%%%%%%%%%%%%%%%%%%%%%%
\subsection{Rugged Mueller potential}
The rugged Mueller potential is a 10D variant of the Mueller potential, a common 2D test problem.
As described in \cite{yuan2024}, the potential is given by

\begin{equation}\label{eq:ruggedmuellerpot}
  V(x) = V_0(x_1, x_2) + \gamma \sin(2 k \pi x_1) \sin(2 k \pi x_2) + \frac{1}{\sigma^2} \sum_{i=3}^{10} x_i^2
\end{equation}
where \(x \in \R^{10}\), and 
\begin{equation}\label{eq:muellerpot}
  V_0(x_1, x_2)  = \sum_{i=1}^4 D_i \exp\left(a_i (x_1 - X_i)^2 + b_i (x_1 - X_i)(x_2-Y_i) + c_i (x_2 - Y_i)^2\right)
\end{equation}
Parameter values are taken from \cite{yuan2024}, and are given in Table~\ref{tab:muellerparams}.
\begin{table}[htbp]
  \centering
  \caption{Parameters of the rugged Mueller potential \eqref{eq:ruggedmuellerpot}--\eqref{eq:muellerpot}.}
  \label{tab:muellerparams}
  \begin{tabular}{c|rrrr}
    \(i\) & 1 & 2 & 3 & 4 \\
    \hline
    \(a_i\) & \(-1\)   & \(-1\)   & \(-6.5\) & \(0.7\) \\
    \(b_i\) & \(0\)    & \(0\)    & \(11\)   & \(0.6\) \\
    \(c_i\) & \(-10\)  & \(-10\)  & \(-6.5\) & \(0.7\) \\
    \(D_i\) & \(-200\) & \(-100\) & \(-170\) & \(15\)  \\
    \(X_i\) & \(1\)    & \(0\)    & \(-0.5\) & \(-1\)  \\
    \(Y_i\) & \(0\)    & \(0.5\)  & \(1.5\)  & \(1\)   \\
    \hline
    \multicolumn{5}{c}{\(\gamma = 9, \quad k = 5, \quad \sigma = 0.05\)}
  \end{tabular}
\end{table}
The metadynamics algorithm additionally requires the gradient of the potential.
Letting 
\begin{equation}
  \begin{aligned}
    \delta X_j &= (x_1 - X_j), \\
    \qquad \delta Y_j &= (x_2 - Y_j),\\
    \qquad S_j &= a_j (\delta X_j)^2 + b_j (\delta X_j)(\delta Y_j) + c_j(\delta Y_j)^2
  \end{aligned}
\end{equation}
the gradient is given by
\begin{equation}\label{eq:grad_mueller_potential}
  \begin{aligned}
    (\nabla V(x))_1 &= 
      \sum_{j=1}^4 \Big( D_j \exp(S_j) \left(2 a_j (\delta X_j) + b_j (\delta Y_j)\right) \Big)
      + \gamma 2 \pi k \cos(2k \pi x_1) \sin(2 k \pi x_2)\\
    (\nabla V(x))_2 &= 
      \sum_{j=1}^4 \Big( D_j \exp(S_j) \left(b_j (\delta X_j) + 2 c_j (\delta Y_j)\right) \Big)
      + \gamma 2 \pi k \sin(2k \pi x_1) \cos(2 k \pi x_2)\\
    (\nabla V(x))_i &= \frac{2}{\sigma^2} x_i \qquad \text{for } i \ge 3
  \end{aligned}
\end{equation}
An efficient vectorized implementation of both the potential and the gradient is given in
\href{https://github.com/jrippeon/amsc665project2/blob/main/src/mueller.py}{\texttt{mueller.py}}.
kjk

%%%%%%%%%%%%%%%%%%%%%%%%%%%%%%%%%%%%%%%%%%%%%%%%%%%%%%%%%%%%%%%%%%%%%%%%%%%%%%%%

\subsection{Euler-Maruyama simulation}
In the metadynamics algorithm (as described in \cite{make_training_data_face}),
the Euler-Maruyama algorithm is used to simulate overdamped Langevin dynamics, given by
\begin{equation}\label{eq:overdamped_langevin}
  dx(t) = - \nabla V(x(t)) \, dt + \sqrt{2 \beta^{-1}} dw(t)
\end{equation}
An Euler-Maruyama step of length \(h\) is given by
\begin{equation}\label{eq:overdamped_langevin_euler_maruyama}
  x_{n+1} = x_n - \nabla V(x_n) h + \sqrt{2 \beta^{-1}} w
\end{equation}
where \(w \sim \mathcal N(0, h)\).
This is implemented in 
\href{https://github.com/jrippeon/amsc665project2/blob/main/src/metadynamics.py}{\texttt{metadynamics.py}},
as \texttt{euler\_maruyama()}.

%%%%%%%%%%%%%%%%%%%%%%%%%%%%%%%%%%%%%%%%%%%%%%%%%%%%%%%%%%%%%%%%%%%%%%%%%%%%%%%%

\subsection{Metadynamics sampling}
\subsubsection{Implementation}
In the metadynamics sampling algorithm,
we aim to sample points more uniformly than we would by simply running the dynamics forward in time.
It consists of a deposition phase, and a sampling phase.

In the deposition phase (\texttt{metadynamics\_deposition()}),
we send an Euler-Maruyama trajectory through the potential landscape,
periodically depositing Guassian ``bumps''.
These fill potential wells, flattening the landscape and allowing more uniform sampling.

These Gaussian ``bumps'' live in a lower-dimensional space of ``collective variables''.
We consider a map \(\xi : \R^d \to \R^m\) to the space of collective variables,
and let 
\begin{equation}
  B_k(x) = h_k \exp \left(-\frac{\abs{\xi(x)- \xi(x_k)}^2}{2 \sigma^2}\right)
\end{equation}
where \(h_k\) is the height of the \(k\)-th bump (constant in our case).

Once the potential landscape is flattened in this way, 
we begin the sampling phase (\texttt{metadynamics\_sampling()}).
To sample, we run more Euler-Maruyama trajectories, 
collecting a point every \(N_\text{skip}\) steps. 
This set of points is then thinned using the \(\delta\)-net algorithm,
which basically consists of iterating through the points in order and at each step deleting all points with distance \(\delta\) of the point under consideration.
The remaining points then constitute the training set.
This is the approach used in \cite{make_training_data_face}.

I introduce one optimization in 
\href{https://github.com/jrippeon/amsc665project2/blob/main/src/metadynamics.py}{\texttt{metadynamics.py}},
and Claude helpfully introduced still greater optimization in 
\href{https://github.com/jrippeon/amsc665project2/blob/main/src/metadynamics_optimized.py}{\texttt{metadynamics\_optimized.py}}.
This brings the runtime from the (estimated, I never actually ran it all the way) 40 minute range down to a minute or so.

My contribution was to recognize that the sampling phase (longest by far) could be modified to take advantage of \texttt{numpy} vectorization,
running a large number of Euler-Maruyama trajectories forward simultaneously.
This doesn't technically match the reference implementation, but it's cheap to implement so I did it.
As long as the total amount of time which passes is long I don't think it will result in large differences in the sample. 

Claude's contribution was to precompute the gradient contribution from the bumps on a uniform grid in collective variable space,
and determine the gradient during sampling by linear interpolation.
This is not too expensive, because the bumps are in the two dimensional collective variable space,
rather than the full 10 dimensional space.
The gradient is given by
\begin{equation}
  -\nabla (V(x_n) + \sum_{j=1}^{N_\text{bumps}}) B_j
\end{equation}
and must be computed at each Euler-Maruyama step.
This change replaces a sum of length \(N_\text{bumps}\), by a single linear interpolation.

An obvious optimization to the \(\delta\)-net algorithm is to use a KD tree.
This is a data structure which partitions \(k\)-dimensional space (in this case the 2-dimensional CV space),
repeatedly in a tree structure.
Nearest neighbor search is then extremely fast.
Claude's implementation uses \texttt{scipy.spacial.cKDTree.query\_ball\_point()},
which is significantly faster than the naive implementation.

\subsubsection{Results}
Sampling was done in 
\href{https://github.com/jrippeon/amsc665project2/blob/main/src/generate_training_set_Yuan.py}{\texttt{generate\_training\_set\_Yuan.py}}.
using the ragged Mueller potential described above,
metadynamics with the parameter values in \autoref{tab:metadynamics_params},
and the \(\delta\)-net to increase sparsity.
This produced 62795 points, the first two coordinates of which are visualized in \autoref{fig:q1_sample_points}.

\begin{table}[htbp]
  \centering
  \begin{tabular}{c |  r | l}
    Parameter & Value & Note\\
    \hline
    \(N_\text{skip}\)  & \(500\) & For deposition\\
    \(N_\text{bumps}\) & \(2000\) \\
    \(h_k\)  & \(5\) & bump height\\
    \(\sigma\) & \(0.05\) \\
    \(h\)  & \(10^{-5}\) & Euler-Maruyama timestep\\
    \(\beta\) & \(0.1\) \\
    \(N_\text{skip}\) & \(100\) & For sampling \\
    \(N_\text{samples}\) & \(2.1 \times 10^{5}\) \\
    \(N_\text{walkers}\) & \(100\)
  \end{tabular}
  \caption{Metadynamics parameters used to generate the training data.
  Largely identical to those used in \cite{yuan2024}.}
  \label{tab:metadynamics_params}
\end{table}

\begin{figure}
  \centering
  \includegraphics[width=0.9\textwidth]{figures/p1_sample_points.png}
  \caption{
    Sample points produced by the metadynamics sampling algorithm,
    projected into the \(x_1\)-\(x_2\) plane.
    Contours of the potential function \(V\) are shown.
  }\label{fig:q1_sample_points}
\end{figure}
\newpage
\printbibliography

\end{document}